\sscp{Austronesian incursion/migration(s): \tsc{\ili{Central} Timor}}{15-06-07}
It is not yet clear where Austronesian-speaking peoples came
from to start what eventually developed into the \tsc{\ili{Central}
Timor} subgroup (see \crf{ch:CT}).
\tsc{\ili{Central} Timor} languages are surrounded by \tsc{Timor-Babar} languages.
There is currently no basis to link these two groups together,
which raises several questions regarding the relationship/connection between them.
There are two logical possibilities:\footnote{
		Our thanks goes
		to Malcolm Ross for outlining these possibilities
		and drawing the problem of the position of \tsc{\ili{Central} Timor}
		to our attention.}

\begin{enumerate}
	\item	\tsc{\ili{Central} Timor} is the remnant of an Austronesian group
				whose extent was curtailed by the later arrival of \tsc{Timor-Babar} languages.
	\item	\tsc{\ili{Central} Timor} contains descendants of a language
				that interpolated itself into the middle of the
				already present \tsc{Timor-Babar} languages.
\end{enumerate}

At present we do not have enough data to decide between these
two scenarios.\footnote{
		\citet[67f]{Edwards2021} discusses some evidence for prehistoric contact between the \tsc{Meto}
		cluster and \tsc{\ili{Central} Timor} languages which may suggests that \tsc{\ili{Central}
		Timor} once extended further west into areas where \ili{Tetun} and/or
		\tsc{Meto} is now spoken,
		though it could also suggest that \tsc{Meto} once extended further
		\ili{east} into areas where \ili{Tetun} is now spoken.}
Another factor that
needs to be considered is the connections/relationships between
\tsc{\ili{Central} Timor}, \tsc{Timor-Babar}, and the neighbouring non-\tsc{An} TAP languages.
More investigation into these connections may by help understand
the pre-history of the \tsc{\ili{Central} Timor} languages.

\citet{SchapperWellfelt2018} discuss contact between the non-\tsc{An}
languages of \ili{Alor} and those of Timor,
including \tsc{\ili{Central} Timor} languages.
They discuss songs sung by some groups in \ili{Alor} which are in a
language unknown to the groups singing but which originate “[{\ldots}]
from Timor and in particular the \ili{Tokodede} ethno-linguistic group [{\ldots}]
The songs were recognized [by \ili{Tokodede} people] although the \ili{Tokodede}
speakers interviewed stated that words in the lyrics had been
changed or corrupted.” Later in the same paper they propose that
the \tsc{\ili{Alor}-Pantar} languages had “[{\ldots}] contact with
Austronesian languages that are no longer extant [{\ldots}]”
(p.112) In the conclusion to their paper they state:

\begin{quote}
		“The evidence indicates that Timor was the dominant player in
		\ili{Alor}-Timor relations and the prestige group from which influence flowed.
		The regions with most significant contact were the \ili{Tokodede} region
		(Likisa-\ili{Maubara}) on Timor and eastern \ili{Alor}, an unsurprising result given the proximity
		of the two islands.
		The contact between the groups on both sides of the Ombai Strait
		appears to have been severed by the creation of international
		borders by colonial powers.” \citep[112]{SchapperWellfelt2018} 
\end{quote}

Our knowledge of the \tsc{\ili{Central} Timor} languages is extremely limited.
For example, \ili{Welaun} was given an ISO code [wlh] only in 2020.
Similarly, \citet[75--93]{Fogaca2017} shows that the \ili{Mambae} [mgm]
cluster has considerable lexical diversity,
with the level of lexical similarity between some varieties dropping
to as low as 63{\%}.
This is negligibly higher than the level of similarity that \citet[35]{Edwards2019}
found between Northwest \ili{Mambae} and \ili{Tokodede}.
More data and careful bottom-up reconstruction may help us better
understand the pre-history of the \tsc{\ili{Central} Timor} languages.